\chapter{Worked Numerical Examples}
\label{app:worked}

This appendix collects fully worked numerical examples that tie the theory to
concrete numbers. Throughout, $\beta=0.0065$, $\Lambda=1.1\times10^{-5}\,$s, and
the six-group $^{235}$U data of Appendix~\ref{app:data} are used unless stated.

\section*{Example E.1: Reactivity conversions}
A reactor has $\keff=1.0013$. The reactivity is
$\rho=(\keff-1)/\keff = 0.0013/1.0013 \approx 1.30\times10^{-3} = 130\,$pcm. In
dollars, $\$=\rho/\beta = 0.0013/0.0065 = 0.20\,\$$. This is a modest,
delayed-critical insertion---well below the one-dollar prompt-critical threshold.

\section*{Example E.2: Prompt jump}
For a step insertion of $\rho=0.20\,\$=0.0013$, the prompt-jump factor is
$\beta/(\beta-\rho)=0.0065/(0.0065-0.0013)=0.0065/0.0052=1.25$. Power jumps
essentially instantly to $1.25\times$ its initial value, then rises slowly on the
delayed-neutron period.

\section*{Example E.3: Stable period from the inhour equation}
For the same $\rho=0.0013$, a one-delayed-group estimate with effective
$\bar\lambda\approx0.08\,$s$^{-1}$ gives a dominant root from
$\rho\approx\beta s/(s+\bar\lambda)$, i.e.\
$s\approx \bar\lambda\,\rho/(\beta-\rho)=0.08\times0.0013/0.0052\approx0.02\,$s$^{-1}$,
so the period $T=1/s\approx 50\,$s. The full six-group inhour solution gives a
similar tens-of-seconds period (cf.\ Figure~\ref{fig:inhour}), confirming a slow,
controllable transient.

\section*{Example E.4: Prompt-critical excursion}
For $\rho=1.5\,\$=0.00975$, the prompt reactivity excess is
$\rho-\beta=0.00325$. The prompt growth rate is
$(\rho-\beta)/\Lambda=0.00325/1.1\times10^{-5}\approx 295\,$s$^{-1}$, an e-folding
time of $\sim3.4\,$ms. In $20\,$ms the power grows by
$e^{295\times0.02}=e^{5.9}\approx 365$. This is the regime in which only prompt
feedback---not control---can intervene.

\section*{Example E.5: Nordheim--Fuchs energy release}
Take $\rho_0=1.2\,\$=0.0078$, so $\rho_0-\beta=0.0013$. Suppose the feedback
parameter is $\gamma=|\alpha|/C$ with $|\alpha|=3\times10^{-5}\,$K$^{-1}$ and
$C=5\times10^{4}\,$J/K, giving $\gamma=6\times10^{-10}\,$J$^{-1}$. The energy at
peak is $E_{\text{peak}}=(\rho_0-\beta)/\gamma=0.0013/6\times10^{-10}\approx
2.2\times10^{6}\,$J, and the total prompt-spike energy is about
$4.3\times10^{6}\,$J. The adiabatic fuel temperature rise is
$\Delta T=E_{\text{total}}/C\approx 87\,$K---a survivable spike, illustrating
self-limitation by negative feedback. Had $\alpha$ been positive, no such cap
would exist.

\section*{Example E.6: Decay-heat heat-up}
A $3000\,$MW$_\text{th}$ core loses cooling just after shutdown with $6\%$ decay
heat, i.e.\ $q=1.8\times10^{8}\,$W. With core thermal mass $C=1.0\times10^{7}\,$J/K,
the adiabatic heat-up rate is $\dot T=q/C=18\,$K/s. To rise $2000\,$K toward fuel
melting takes about $110\,$s---under two minutes---quantifying the absolute need
for assured decay-heat removal.

\section*{Example E.7: Void feedback loop gain (illustrative)}
Suppose a power increase of $10\%$ raises the core-average void fraction by
$0.02$, and the void coefficient is $\alpha_v=+1.5\,\$$ per unit void fraction
(an RBMK-like positive value at low power). The feedback reactivity is
$\Delta\rho=+0.02\times1.5=+0.03\,\$$. Positive feedback of this sign and size,
unopposed by sufficient prompt Doppler, drives the power further up---the unstable
loop. The same calculation with $\alpha_v<0$ (a PWR) yields negative feedback and
a self-correcting response.
